% ECONOMICS_PROBLEM: {"id": "D63-3", "title": "Complete EFX allocations for additive goods", "jel_code": "D63", "source_id": "OP-0186"}
% FIRST_POSTED: 2026-10-09
% ATTEMPT_STATUS: unresolved
% ENTRY_KIND: research_attempt
% !TeX program = pdflatex
\documentclass[11pt,a4paper]{article}
\usepackage[T1]{fontenc}
\usepackage[utf8]{inputenc}
\usepackage{lmodern}
\usepackage[margin=24mm,headheight=15pt]{geometry}
\usepackage{amsmath,amssymb,amsthm,mathtools}
\usepackage{booktabs,longtable,array,enumitem}
\usepackage{microtype}
\usepackage{xurl}
\usepackage[hidelinks]{hyperref}
\usepackage{fancyhdr}
\newcommand{\E}{\mathbb E}
\newcommand{\R}{\mathbb R}
\newcommand{\N}{\mathbb N}
\newcommand{\ind}{\mathbf 1}
\DeclareMathOperator{\Reg}{Reg}
\DeclareMathOperator{\supp}{supp}
\newtheorem{theorem}{Theorem}[section]
\newtheorem{lemma}[theorem]{Lemma}
\newtheorem{proposition}[theorem]{Proposition}
\newtheorem{corollary}[theorem]{Corollary}
\theoremstyle{definition}
\newtheorem{definition}[theorem]{Definition}
\theoremstyle{remark}
\newtheorem{remark}[theorem]{Remark}
\setlength{\parindent}{0pt}
\setlength{\parskip}{4pt plus 1pt}
\setlength{\emergencystretch}{2em}
\setcounter{tocdepth}{1}
\allowdisplaybreaks[2]
\pagestyle{fancy}
\fancyhf{}
\fancyhead[R]{\small Open Problems Econ}
\fancyfoot[C]{\thepage}
\renewcommand{\headrulewidth}{0.3pt}
\hypersetup{pdfauthor={GPT 6 Astra Pro},pdfsubject={Checked partial results and explicit remaining gaps}}

\fancyhead[L]{\small\texttt{D63-3} -- Research attempt}
\title{Research attempt: \texttt{D63-3}\\[5pt]\large Complete EFX for additive goods}
\author{GPT 6 Astra Pro}
\date{9 October 2026}
\begin{document}
\maketitle
\noindent\textbf{Scope of this report.} The main target remains UNRESOLVED. Fully proved subclaims and counterexamples are identified locally. Computational checks reported here are supplementary to the proofs. Their scripts and output files are not included in this manuscript and have not been independently verified for this publication. No novelty or priority claim is made.\par
\section{FORMAL RESTATEMENT}
Let $N=\{1,\ldots,n\}$ with $n\geq1$, let $G$ be a finite set of $m\geq0$ goods, and let each $v_i:G\to\mathbb R_{\geq0}$ extend additively by $v_i(S)=\sum_{g\in S}v_i(g)$. A complete allocation is an ordered partition $(A_1,\ldots,A_n)$ of $G$, with empty bundles permitted. Define
\begin{align*}
 \mathrm{EFX}_0:\quad &\forall i,j\ \forall g\in A_j,
             \quad v_i(A_i)\geq v_i(A_j\setminus\{g\}),\\
 \mathrm{EFX}_+:\quad &\forall i,j\ \forall g\in A_j\text{ with }v_i(g)>0,
             \quad v_i(A_i)\geq v_i(A_j\setminus\{g\}).
\end{align*}
The two universal existence assertions are
\[
 \forall n\geq1\ \forall m\geq0\ \forall(v_i(g))\in\mathbb R_+^{n\times m}
 \ \exists\text{ complete allocation satisfying the selected EFX convention}.
\]
No efficiency, welfare optimality, or polynomial-time construction is required. An empty compared bundle has no goods to test. For $n=1$ or $m=0$, every allocation is EFX. Goods are not chores: negative values are excluded.

\paragraph{Literal statement and hygiene.}
The problem statement contains overlapping formulations; they are treated as one additive-goods existence question, retaining the zero-value convention as a distinct definition. The conventions are not identical for a fixed allocation. Nevertheless, we prove their universal existence questions equivalent. Partial allocations, approximate EFX, and nonadditive counterexamples cannot resolve this complete additive target.

\section{QUICK LITERATURE/CONTEXT CHECK}
Chaudhury, Garg and Mehlhorn prove complete EFX existence for three additive agents \cite{d63-three}. A 9 August 2026 preprint by Alkassar, Fouz and Mehlhorn reports complete strong $\mathrm{EFX}_0$ for four nonnegative additive agents and at most nine goods, using a certificate corpus \cite{d63-four}. The present computation does not rerun or certify that external corpus. These results delimit small counterexample searches; they do not supply a general-agent proof here. The constructive subclasses below are proved directly. Sources checked through 9 October 2026.

\section{ATTACK PLAN}
\paragraph{Proof strategies.}
First, remove ambiguity at zero values by perturbation. Second, seek allocation rules with a monotone potential: least-loaded greedy allocation for common values and leximin improvement for binary values. Third, test whether the binary exchange argument survives heterogeneous weights.

\paragraph{Disproof strategies.}
Exhaust small integer valuation arrays, stress near-identical and highly unbalanced valuations, and formulate a finite rational certificate excluding every complete allocation. The heterogeneous leximin strategy fails on an explicit small instance, but that instance still admits EFX.

\paragraph{Landscape and tools.}
This is finite combinatorial fair division. Positive perturbation handles zero tests; finiteness of allocation space permits subsequence selection; decreasing-order greedy allocation compares last-item sizes; cut-and-choose solves two agents; leximin potentials justify binary transfers; exact enumeration checks small arrays; and rational linear feasibility describes a finite disproof certificate.

\section{WORK}
\subsection{Zero conventions and rational counterexample reduction}
\begin{theorem}[Equivalence of the two universal existence assertions]
For any fixed pair $(n,m)$, every nonnegative additive instance admits $\mathrm{EFX}_+$ if and only if every such instance admits $\mathrm{EFX}_0$.
\end{theorem}
\begin{proof}
The reverse implication is immediate from inclusion of the tested goods. For the forward direction, fix an arbitrary nonnegative valuation array $v$ and put $v_i^{(k)}(g)=v_i(g)+1/k$. All item values are strictly positive. By the assumed universal $\mathrm{EFX}_+$ assertion, there is a complete allocation $A^{(k)}$ satisfying all the $\mathrm{EFX}_0$ inequalities for $v^{(k)}$. There are only $n^m$ complete allocations, so one allocation $A$ occurs for infinitely many $k$. Along that subsequence, each of its finitely many weak linear inequalities passes to the limit. They are exactly the $\mathrm{EFX}_0$ inequalities for $v$.

This is an equivalence of universal assertions, not a claim that an arbitrary $\mathrm{EFX}_+$ allocation at zero values is $\mathrm{EFX}_0$. For example, agent 1 with an empty bundle may value the two goods in agent 2's bundle as 1 and 0. Removing the valued good leaves zero and passes the positive test; removing the zero-valued good leaves 1 and fails the zero-tolerant test. Give agent 2 value zero for everything to make the other tests vacuous or satisfied.
\end{proof}

\begin{lemma}[Strictly positive rational witnesses suffice]
If some real nonnegative valuation array has no complete $\mathrm{EFX}_0$ allocation, then there is a strictly positive rational valuation array of the same dimensions with no such allocation. It can be scaled to a positive integer array.
\end{lemma}
\begin{proof}
The claim is vacuous when $m=0$. Otherwise, for each of the finitely many allocations $A$, choose a violated inequality, with gap
\[
 d_A=v_i(A_j\setminus\{g\})-v_i(A_i)>0.
\]
Let $d=\min_A d_A>0$. Here $i\ne j$, because nonnegativity rules out a self-comparison violation. The two bundles are disjoint, so the linear gap uses at most $m$ coefficients, each of absolute value one. Choose a strictly positive rational value within $d/(2m)$ of each original entry; density of positive rationals in $[0,\infty)$ permits this even at original zeros. Every selected gap changes by less than $d/2$ and remains positive, excluding every allocation. Multiply all entries by a common positive denominator to get a positive integer array without changing the EFX inequalities.
\end{proof}

\paragraph{An exact disproof search, not merely a floating-point test.}
For fixed $n,m$, enumerate all choices assigning to each allocation $A$ a potential witness triple $(i,j,g\in A_j)$. For each such finite choice solve the rational linear system
\[
 v_i(g)\geq1\quad\text{for all entries},\qquad
 v_i(A_j\setminus\{g\})-v_i(A_i)\geq1
       \quad\text{for every selected allocation witness}.
\]
Feasibility gives an explicit positive rational counterexample. Conversely, any strict rational counterexample can be multiplied by a sufficiently large positive constant to make every entry and selected positive gap at least one, so some enumerated system is feasible. Rational linear programming is an exact decision procedure for each system. This is extremely large but finite for fixed dimensions; enumeration over $(n,m)$ is a semidecision procedure for falsity of the universal assertion. It is not an efficient construction or a proof that a feasible witness system exists.

\subsection{Four fully proved existence classes}
\begin{lemma}[Identical valuations]
Every instance in which all agents have the same additive valuation admits a complete $\mathrm{EFX}_0$ allocation.
\end{lemma}
\begin{proof}
Sort the goods in nonincreasing common value and assign them one at a time to a bundle of currently minimum common value. Fix any final nonempty bundle $A_j$ and let $g_*$ be its last assigned good. Immediately before that assignment, the common value of every bundle $A_i$ then present was at least the value of $A_j$ then present, namely $v(A_j)-v(g_*)$. Each bundle's value can only increase afterward, so final $v(A_i)\geq v(A_j)-v(g_*)$. By the decreasing processing order, $g_*$ has minimum value among goods in $A_j$. Thus for every $g\in A_j$,
$v(A_i)\geq v(A_j)-v(g_*)\geq v(A_j)-v(g)=v(A_j\setminus\{g\})$.
Zero-valued goods and ties obey the same weak inequalities. Empty bundles $A_j$ require no test.
\end{proof}

\begin{corollary}[Two agents]
Every two-agent nonnegative additive instance has a complete $\mathrm{EFX}_0$ allocation.
\end{corollary}
\begin{proof}
Agent 1 partitions the goods into two bundles using the identical-valuation lemma with two copies of its own valuation. Agent 2 takes a bundle of maximum value under $v_2$, leaving the other to agent 1. Agent 2 does not envy even before removal, and nonnegativity preserves its inequality after any removal. Agent 1 satisfies the EFX inequalities regardless of which of the two bundles it receives, because the construction established both directions under $v_1$.
\end{proof}

\begin{lemma}[At most one good beyond the number of agents]
Every instance with $m\leq n+1$ admits complete $\mathrm{EFX}_0$.
\end{lemma}
\begin{proof}
If $m\leq n$, give at most one good to each agent. Every compared nonempty bundle becomes empty after removal, so all inequalities hold. If $m=n+1$, let the first $n-1$ agents successively choose a favorite remaining good, then give the last two goods to agent $n$. Any earlier agent values its chosen good at least as much as each of those last two goods, since they were available at its turn. Removing either good from the last bundle leaves only one, so its EFX comparison holds. All other compared bundles are singletons, and removing their good leaves zero. The last agent also compares only with singletons. When $n=1$, this construction gives its sole agent both goods and the inequalities follow from nonnegativity.
\end{proof}

\begin{theorem}[Binary and agent-scaled binary valuations]
Every additive instance with $v_i(g)\in\{0,1\}$ has a complete $\mathrm{EFX}_0$ allocation. The same conclusion holds when, for each agent $i$, values lie in $\{0,w_i\}$ for an arbitrary $w_i\geq0$.
\end{theorem}
\begin{proof}
First remove goods valued zero by everyone. Each remaining good is valued 1 by at least one agent. Among assignments giving each such good to an agent who values it 1, choose one lexicographically maximizing the vector of own utilities sorted in nondecreasing order. This maximum exists because there are finitely many assignments. In these nonwasteful assignments, each owner's utility is its bundle's cardinality.

Suppose this allocation fails $\mathrm{EFX}_0$. Then for some $i,j,g\in A_j$,
$u_i=v_i(A_i)=a<v_i(A_j\setminus\{g\})$.
Integrality gives $v_i(A_j\setminus\{g\})\geq a+1$. Hence
$u_j=|A_j|\geq |A_j\setminus\{g\}|+1\geq a+2$.
There is a good $h\in A_j$ valued 1 by $i$. Move it from $j$ to $i$. The assignment stays nonwasteful, and its two changed utilities go from $(a,b)$ to $(a+1,b-1)$ with $b\geq a+2$. This strictly improves the increasingly sorted utility vector in lexicographic order: all entries below $a$ remain unchanged; the number of entries equal to $a$ decreases by one; and both replacement entries exceed $a$. At the first changed sorted position the new entry is therefore larger. This contradicts maximality, so the allocation is $\mathrm{EFX}_0$.

Now assign all globally zero goods to an agent $j_0$ of minimum own utility in that allocation. For every $i$,
$u_i\geq u_{j_0}=|A_{j_0}|\geq v_i(A_{j_0})$
before the zero goods are added. Thus removing any newly added zero good from that bundle cannot create envy; it leaves the same bundle value as before addition, which is at most $u_i$. Removing an old good produces its old EFX inequality. All other comparisons and own utilities are unchanged. This completes the allocation while preserving $\mathrm{EFX}_0$.

For agent-scaled binary values, divide each positive-valued agent's entire valuation by $w_i$. This positive scaling does not change any of that agent's EFX inequalities. Agents with $w_i=0$ already have only zero values. Apply the binary case and then undo the scaling.
\end{proof}

\subsection{A failed extension, with a smallest transparent witness}
Take two agents and four goods, with agent 1 valuing every good at 100 and agent 2 valuing every good at 1. An allocation giving $k$ goods to agent 1 has utility pair $(100k,4-k)$. The increasingly sorted lexicographically maximal utility pair is $(3,100)$, obtained by giving one good to agent 1 and three to agent 2. Agent 1 then values its own bundle at 100 but the other bundle after any removal at 200. Every such leximin allocation fails EFX. Giving two goods to each agent is EFX, so this is not a disproof of the main conjecture. It pinpoints why the binary proof cannot be extended by retaining raw heterogeneous utilities: a transfer no longer changes the two owners' utilities by equal units.

\section{VERIFICATION}
The reported exact enumeration tests all 729 arrays in $\{0,1,2\}^{2\times3}$ and all 4096 arrays in $\{0,1\}^{3\times4}$. Every instance has an $\mathrm{EFX}_0$ allocation. For the binary instances it also checks the leximin-with-zero-goods construction, with no failure. The computation additionally tests 128 fixed-seed random and near-identical four-agent, seven-good arrays, examining all $4^7=16{,}384$ allocations per array. Each has an EFX allocation; the observed counts range from 30 to 630. These are finite checks, not evidence sufficient for arbitrary real valuations or arbitrary dimensions. The fixed seed is 20261009.

Pseudocode is: enumerate a valuation array; enumerate every assignment of each good to one owner; compute each agent's value for each bundle; reject an allocation if any tested removal leaves strict envy; record surviving allocations. All these calculations use integers. A second loop selects a nonwasteful lexicographic maximum in binary instances, restores globally zero goods as proved, and checks every inequality again.

The perturbation proofs use only finitely many allocations and finite strict margins. They do not require a continuous selection of allocations. The binary proof explicitly handles removing a zero-valued good, the issue that would defeat a positive-only argument. The rational-certificate search quantifies over a witness for every complete allocation; finding one bad allocation is not enough. The external nine-good certificate corpus has not been independently executed in this package.

\section{FINAL}
\textbf{LABEL: UNRESOLVED} for general complete additive EFX existence.

\paragraph{Strongest checked partial results.}
The two universal zero-value conventions are equivalent. Complete $\mathrm{EFX}_0$ exists for identical valuations, two agents, $m\leq n+1$, and arbitrary numbers of agents with agent-scaled binary values. Any genuine counterexample can be chosen with strictly positive rational, hence positive integer, valuations. A finite linear witness system exactly captures a disproof at fixed dimensions.

\paragraph{Exact first gap.}
No allocation potential or exchange operation is proved to yield complete EFX for arbitrary heterogeneous additive weights, and no rational witness system excluding every allocation is exhibited. In particular, the unit-transfer improvement lemma used for binary valuations has no proved weighted replacement.

\paragraph{Three next moves.}
(1) Search exact witness systems in the smallest dimensions outside established positive results, using symmetry reduction but preserving all allocation witnesses. (2) Develop a normalized or multi-coordinate potential whose every EFX violation admits a strictly improving feasible exchange. (3) Study near-identical positive integer arrays with controlled perturbations, recording which EFX allocations disappear as inequalities cross zero.

\paragraph{Minimal counterexample structure.}
The direct proofs exclude two agents, common values, agent-scaled binary values, and $m\leq n+1$. The published three-agent result further requires at least four agents. Accepting the reported 2026 four-agent theorem, a four-agent counterexample would need at least ten goods. These literature exclusions are not being re-certified by the smaller searches here. A minimal witness can have all item values strictly positive and rational.

\begin{thebibliography}{9}
\bibitem{d63-three} B. R. Chaudhury, J. Garg, and K. Mehlhorn. \emph{EFX Exists for Three Agents}. arXiv:2002.05119; full version of EC 2020. \url{https://arxiv.org/abs/2002.05119}.
\bibitem{d63-four} E. Alkassar, M. Fouz, and K. Mehlhorn. \emph{Complete EFX Allocations Exist for Four Additive Agents and Up to Nine Goods}. arXiv:2608.08590, August 2026. \url{https://arxiv.org/abs/2608.08590}.
\end{thebibliography}

\end{document}
